% RESEARCH_DRAFT_SHA256: ad4c06ca0060d14e44dff86d0b51e494953e11d36e1dd0402483833fe521ccbb
% Layout and source-reference wrappers expanded for the web renderer.
\documentclass[11pt]{article}
\usepackage[margin=25mm]{geometry}
\usepackage{fontspec,unicode-math}
\setmainfont{Latin Modern Roman}
\setmathfont{Latin Modern Math}
\usepackage{amsmath,amssymb,mathtools,xurl,hyperref}
\hypersetup{colorlinks=true,urlcolor=blue,linkcolor=blue}
\setlength{\parindent}{0pt}
\setlength{\parskip}{5pt}
\setlength{\emergencystretch}{3em}
\newcommand{\sref}[2]{\hyperlink{source:#1:#2}{[S#2]}}
\newcommand{\cataloguescope}[1]{\paragraph{Source question.}#1}
\newcommand{\floor}[1]{\left\lfloor#1\right\rfloor}
\begin{document}

\section*{A higher-power witness from one nonresidue: exact boundary search and root extraction}
\noindent catalogue entry; catalogue entry. Research attempt dated 11 October 2026.
\paragraph{Result in scope.}This manuscript contains a source-specific attempt and complete proofs of the stated auxiliary calculations. It makes no new solution or novelty claim for the entire target.
\section{Original problem and preserved definition}
The following definition and source list are reproduced from the repository. Any earlier inspection dates and reported literature claims in this preserved text describe that source record, not additional verification by this attempt.
% BEGIN PRESERVED DEFINITION
\subsection{Definitions and mathematical statement}

\paragraph{Shared notation.}Standard number systems are denoted by $\mathbb N,\mathbb Z,\mathbb Q,\mathbb R,\mathbb C$; index ranges, computational models and specialized notation follow the cited problem source and the definitions below.


Polynomial time means a bound polynomial in the finite input length; for an integer input, its binary length is proportional to its logarithm. Write $p-1=2^\ell m$ with m odd and $\ell\ge3$. A $2^{\ell-1}$th power in $\mathbb F_p$ is an element $y^{2^{\ell-1}}$ for some $y\in\mathbb F_p$. The target is to produce x for which $x^{2^{\ell-1}}+1$ is not such a power, in deterministic time polynomial in the binary length of p.
\paragraph{Statement recorded in the supplied source.}
Question 12. Let p be prime, write p − 1 = 2 `m where m is odd, and assume

` ≥ 3. Find in deterministic polynomial time x ∈ Fp such that x2`−1

+ 1 is not a 2 `−1th power. (Kedlaya) \href{https://aimath.org/WWN/primesinp/primesinp.pdf}{[S2]}
\paragraph{Scope qualification.}
The superscript nesting and the symbol ell have been flattened by extraction; the definitions use the exponents printed in the primary notes.
\paragraph{Source status.}
The supplied catalogue labels this record partially solved. Its source-status metadata was checked against the supplied snapshot on 2026-10-11. This import does not independently establish current open status or mathematical correctness.
\subsection{Short English statement}
Can a specified higher-power nonresidue expression be found deterministically in polynomial time?
\subsection{Sources}
\begin{itemize}
\item[S1] Ulam AI and the UnsolvedMath Contributors, UnsolvedMath: A Curated Collection of Open Mathematics Problems, catalogue entry. Catalogue curation and metadata: CC BY 4.0. The supplied source question is retained; definitions and the English summary are editorial additions.
\url{https://huggingface.co/datasets/ulamai/UnsolvedMath}.
\item[S2] AIM Workshop Problem Lists, original problem source linked by the catalogue
\url{https://aimath.org/WWN/primesinp/primesinp.pdf}.
\end{itemize}
% END PRESERVED DEFINITION
\subsection{Verbatim supplied catalogue question}
Question 12. Let p be prime, write p − 1 = 2 `m where m is odd, and assume

` ≥ 3. Find in deterministic polynomial time x ∈ Fp such that x2`−1

+ 1 is not a 2 `−1th power. (Kedlaya)

\section{Attempt and supplied-data condition}
Write $p-1=2^\ell m=2km$, with $m$ odd, $\ell\ge3$ and
$k=2^{\ell-1}$. The target is $x$ for which $x^k+1$ is not a
$k$th power, including zero among the powers. We prove a deterministic
construction if one quadratic nonresidue $z$ is supplied. Producing
that supplied nonresidue for arbitrary $p$ remains a separate step,
so the conditional construction is not the original unconditional
algorithm.

\paragraph{The power subgroup and its membership test.}
The multiplicative group of a finite field is cyclic. Here is the
needed elementary proof. In any finite abelian group one can construct
an element whose order is the exponent (the least common multiple of
all element orders): for each prime choose an element containing the
maximal prime-power order, raise it to remove its other primary parts,
and multiply these commuting elements of coprime orders. In a finite
subgroup of a field, if this exponent is $e$, all group elements are
roots of $X^e-1$, so their number is at most $e$. The constructed
element gives $e$ distinct elements, hence it generates the group.
Thus
$H=(\mathbb F_p^*)^k$ has size $2m$, equals the roots of
$u^{2m}=1$, and contains $-1$. Each $h\in H$ has exactly $k$
preimages under $x\mapsto x^k$, as is immediate from exponents of
a generator. Let $H_0=H\cup\{0\}$. Membership in $H_0$ is tested
by equality to zero or by the displayed power test. Since $k$ is even,
$H$ consists of squares, and a quadratic nonresidue is outside $H_0$.

\paragraph{Nonmonotone bisection finds an additive boundary.}
Represent the supplied nonresidue $z$ as an integer in $[1,p-1]$.
Maintain integers $a<b$ with $a\in H_0$ and $b\notin H_0$, starting
at $a=0,b=z$. Test the integer midpoint: if it is in $H_0$, replace
$a$ by it; otherwise replace $b$. Each step preserves the opposite
endpoint labels and halves the interval length up to rounding, so after
$O(\log p)$ tests we have $b=a+1$. No monotonicity of the subgroup
labels is assumed. Since one belongs to $H_0$, the final $a$ cannot
be zero. Set $h=a\in H$; then $h+1\notin H_0$.

\paragraph{An explicit $k$th root of the boundary value.}
If $m>1$, choose an integer $e$ with $ke\equiv1\pmod m$ by the
integer Euclidean algorithm; this is possible since $k$ is a power of
two and $m$ is odd. For $m=1$ set $e=0$. In either case
$ke=1+tm$ with $t$ odd, since $ke$ is even and $m$ is odd.
For $h\in H$, $h^m$ equals $1$ or $-1$, so
\[
(h^e)^k=h(h^m)^t=h h^m.
\]
A nonresidue $z$ has $z^{(p-1)/2}=-1$, by the cyclic-group
criterion for squares. Therefore $c=z^m$ satisfies $c^k=-1$.
If $h^m=1$, take $x=h^e$; if $h^m=-1$, take $x=h^e c$.
The displayed equations prove $x^k=h$ in both cases. Consequently
$x^k+1=h+1$ is not a $k$th power, as required.
The membership tests and exponentiations use exponents with
$O(\log p)$ bits; binary powering uses $O(\log p)$ field
multiplications per test. Together with the Euclidean algorithm and
$O(\log p)$ bisections this is polynomial in the binary length of
$p$, conditional on the supplied nonresidue.

\paragraph{Existence alone does not locate a witness efficiently.}
The proper set $H_0$ contains zero and one. Taking the least positive
integer outside it gives an additive boundary as above, so such a
boundary exists without any algorithmic assumption. Each boundary
$h\in H$ has $k$ distinct $k$th roots, giving at least $k$ witnesses.
This bound can be exact: for $p=17$, $m=1$, $k=8$ and
$H=\{1,-1\}$. Only $h=1$ has $h+1$ outside $H_0$, since
$1+1=2\notin\{0,1,-1\}$ while $-1+1=0$; $x=0$ also fails.
There are exactly eight witnesses. An existence proof or a brute-force
scan over $p$ values is not a deterministic bound polynomial in $\log p$.

\section{Audit and gap}
The primary Question 12's nested exponents and the nearby Agrawal
nonresidue discussion were checked. Exact Python tests on every prime
$p<200$ with $8\mid p-1$, and every supplied quadratic nonresidue for
each such prime, checked the bisection, extracted root and witness
condition; the $p=17$ witness count was eight. These finite tests only
verify implementation examples of the proved conditional construction.
The audit checked zero as a power, the $m=1$ inverse convention,
nonmonotone bisection and the supplied-nonresidue dependency. No
unconditional fast nonresidue routine is proved, and no new resolution
of the source's deterministic all-primes target is claimed.

\section{Process and contributions}
Mathematics and model-assisted adversarial self-audit: OpenAI gpt-6-astra, ultra reasoning, as dispatched for this research batch; exact serving revision not exposed. Codex (GPT-6; exact serving revision not exposed) materially assisted with orchestration and format review. No human mathematical authorship, mathematical edits, or independent human proof checking is asserted.
Contribution prompt A was applied to the full definition and supplied source row, then prompt B was applied adversarially to the derived results. The item-specific audit above records the scope checks and rejected extrapolations. Local assembly and structural checks do not establish mathematical correctness. The source catalogue status is preserved in TOP\_PROBLEM and is separate from this attempt status.
\paragraph{Integration review.} Codex (GPT-6; exact serving revision not exposed) reviewed the written arguments and expanded layout and source-reference wrappers. Repository metadata and the production web renderer were checked using Python 3.13 and JavaScript with MathJax 3.2.2. This model-assisted review is not independent human proof verification; no PDF compilation or formal proof checking is asserted.

\section{Completion Estimate}
\noindent\textbf{COMPLETION: 0\%}
The proved material is a rederivation of known or auxiliary facts rather than a substantiated new advance on the original target. A complete deterministic reduction from one supplied quadratic nonresidue is proved, but producing that input unconditionally for arbitrary p within the required time remains missing.
\section{References}
The complete original Sources subsection above is retained, including catalogue attribution and licence notices. The additional list records primary sources actually consulted for this attempt on 11 October 2026; uncited source links reproduced above were not all rechecked.
\begin{itemize}
\item American Institute of Mathematics, Future directions in algorithmic number theory, workshop 24--28 March 2003, report version 30 April 2003; original question, relevant lecture context and adjoining remarks inspected on 11 October 2026. \url{https://aimath.org/WWN/primesinp/primesinp.pdf}.
\end{itemize}
\section{Final status}
\textbf{Unresolved in this attempt.} A complete deterministic reduction from one supplied quadratic nonresidue is proved, but producing that input unconditionally for arbitrary p within the required time remains missing.
\end{document}
